\documentclass[a4paper]{article}

\input{style/ch_xelatex.tex}
\input{style/scala.tex}

\graphicspath{ {images/} }
\usepackage{ctex}
\usepackage{graphicx}
\usepackage{color,framed}%文本框
\usepackage{listings}
\usepackage{caption}

\usepackage{hyperref}
\hypersetup{hidelinks,
	colorlinks=true,
	allcolors=black,
	pdfstartview=Fit,
	breaklinks=true}


\usepackage{amssymb}
\usepackage{enumerate}
\usepackage{xcolor}
\usepackage{bm} 
\usepackage{lastpage}%获得总页数
\usepackage{fancyhdr}
\usepackage{tabularx}  
\usepackage{geometry}
\usepackage{minted}
\usepackage{graphics}
\usepackage{subfigure}
\usepackage{float}
\usepackage{pdfpages}
\usepackage{pgfplots}
\pgfplotsset{width=10cm,compat=1.9}
\usepackage{multirow}
\usepackage{footnote}
\usepackage{booktabs}
\usepackage{listings}
\usepackage{seqsplit}
\usepackage[table]{xcolor}
\usepackage{array}
%-----------------------伪代码------------------
\usepackage{algorithm}  
\usepackage{algorithmicx}  
\usepackage{algpseudocode}  
\floatname{algorithm}{Algorithm}  
\renewcommand{\algorithmicrequire}{\textbf{Input:}}  
\renewcommand{\algorithmicensure}{\textbf{Output:}} 
\usepackage{lipsum}  
\makeatletter
\newenvironment{breakablealgorithm}
  {% \begin{breakablealgorithm}
  \begin{center}
     \refstepcounter{algorithm}% New algorithm
     \hrule height.8pt depth0pt \kern2pt% \@fs@pre for \@fs@ruled
     \renewcommand{\caption}[2][\relax]{% Make a new \caption
      {\raggedright\textbf{\ALG@name~\thealgorithm} ##2\par}%
      \ifx\relax##1\relax % #1 is \relax
         \addcontentsline{loa}{algorithm}{\protect\numberline{\thealgorithm}##2}%
      \else % #1 is not \relax
         \addcontentsline{loa}{algorithm}{\protect\numberline{\thealgorithm}##1}%
      \fi
      \kern2pt\hrule\kern2pt
     }
  }{% \end{breakablealgorithm}
     \kern2pt\hrule\relax% \@fs@post for \@fs@ruled
  \end{center}
  }
\makeatother
%------------------------代码-------------------
\RequirePackage{listings}
\RequirePackage{xcolor}
\definecolor{dkgreen}{rgb}{0,0.6,0}
\definecolor{gray}{rgb}{0.5,0.5,0.5}
\definecolor{mauve}{rgb}{0.58,0,0.82}
\lstset{
	numbers=left,  
	frame=tb,
	aboveskip=3mm,
	belowskip=3mm,
	showstringspaces=false,
	columns=flexible,
	framerule=1pt,
	rulecolor=\color{gray!35},
	backgroundcolor=\color{gray!5},
	basicstyle={\ttfamily},
	numberstyle=\tiny\color{gray},
	keywordstyle=\color{blue},
	commentstyle=\color{dkgreen},
	stringstyle=\color{mauve},
	breaklines=true,
	breakatwhitespace=true,
	tabsize=3,
}

%-------------------------页面边距--------------
\geometry{a4paper,left=2.3cm,right=2.3cm,top=2.7cm,bottom=2.7cm}
%-------------------------页眉页脚--------------
\usepackage{fancyhdr}
\pagestyle{fancy}
\lhead{\kaishu 数据安全}% @这里是填写什么名字的实验报告
\chead{零知识证明实践} % 这里是填写本次实验的名字
\rhead{2312966 {\CJKfontspec{simkai.ttf} 林晖鹏}}%加粗\bfseries 
\lfoot{}
\cfoot{\thepage}
\rfoot{}
\renewcommand{\headrulewidth}{0.1pt}  
\renewcommand{\footrulewidth}{0pt}%去掉横线
\newcommand{\HRule}{\rule{\linewidth}{0.5mm}}%标题横线
\newcommand{\HRulegrossa}{\rule{\linewidth}{1.2mm}}
\setlength{\textfloatsep}{10mm}%设置图片的前后间距
\newcolumntype{Y}{>{\centering\arraybackslash\ttfamily}X}
\newcommand{\codewrap}[1]{\texttt{\seqsplit{#1}}}
%--------------------新增：字体与章节风格配置（匹配目标样式）--------------------
% % 1. 中文字体全局设置：标题黑体、正文宋体
% \setCJKmainfont{PingFang}[AutoFakeBold=true]   % 正文宋体（Windows）| macOS替换为Songti SC | Linux替换为Noto Serif CJK SC
% \setCJKsansfont{Heiti SC}[AutoFakeBold=true]   % 标题黑体（Windows）| macOS替换为Heiti SC | Linux替换为Noto Sans CJK SC
% \setCJKmonofont{Heiti SC}[AutoFakeBold=true]


% 3. 正文格式优化：首行缩进2字符、行高舒展
\usepackage{indentfirst}  % 章节后首段也缩进
\setlength{\parindent}{2em}  % 首行缩进2个中文字符
\linespread{1.5}  % 行高1.3倍（更易读）

%--------------------新增结束--------------------
%--------------------文档内容--------------------

\begin{document}
\renewcommand{\contentsname}{目\ 录}
\renewcommand{\appendixname}{附录}
\renewcommand{\appendixpagename}{附录}
\renewcommand{\refname}{参考文献} 
\renewcommand{\figurename}{图}
\renewcommand{\tablename}{表}
\renewcommand{\today}{\number\year 年 \number\month 月 \number\day 日}

%-------------------------封面----------------

\renewcommand {\thefigure}{\thesection{}.\arabic{figure}}%图片按章标号
\renewcommand{\figurename}{图}
\renewcommand{\contentsname}{目录}  
\cfoot{\thepage\ of \pageref{LastPage}}%当前页 of 总页数

\renewcommand{\abstractname}{\textbf{\Large 摘要}} % 调整摘要标题的字体大小

\newpage
\begin{center}
    \huge{\textbf{《数据安全》实验报告}}\\
    \vspace {0.5em}
    \Large {\textbf {零知识证明实践}}% 上面还有一个地方需要修改
\end{center}


\begin{center}
    \textbf{姓名}：\underline{林晖鹏} \quad
    \textbf{学号}：\underline{2312966} \quad
    \textbf{专业}：\underline{计算机科学与技术}
\end{center}

\vspace{1em} % 添加垂直间距

\tableofcontents

% \vspace{2em} % 添加垂直间距
% % 摘要部分

% \noindent{\Large\textbf{摘要}}
% 本次实验围绕 Paillier 加法同态加密展开，使用 Python 的 \texttt{phe} 库实现了一个基础版 PIR（Private Information Retrieval，隐私信息获取）模拟系统。实验中，客户端先生成 Paillier 公私钥，再构造加密后的选择向量，服务器在不知道查询位置的情况下完成密文域计算，最后由客户端解密得到目标消息。在此基础上，进一步完成了扩展实验：客户端保存 AES-256-GCM 对称密钥，服务器仅保存经过对称加密后的密文，客户端通过多轮 PIR 取回目标密文后再在本地解密出原文。实验结果表明，基础版与扩展版流程都可以正确运行，能够较好体现 Paillier 同态加密在隐私查询场景中的应用方式。

% \vspace{2em}
% \noindent % 取消首行缩进 

% \vspace{1em} % 添加垂直间距
% \noindent % 顶格
% \textbf{关键字：} Paillier加密；同态加密；隐私信息获取；AES-GCM；\texttt{phe}库

\newpage

%--------------------------Title--------------------------------

\section{实验要求}
参考教材实验3.1，假设Alice希望证明自己知道如下方程的解$ x^3 + x + 5 = out $，其中out是大家都知道的一个数，这里假设out为35，而$x=3$就是方程的解请实现代码完成证明生成和证明的验证。

%---------------------------------------------------------------
\section{实验原理}
\subsection{实验背景}
在传统认证或秘密验证中，证明者往往需要直接提交密码、密钥或者答案本身，验证者才能确认其是否掌握秘密。这种方式虽然简单，但也存在明显风险：一旦秘密在传输或存储过程中泄露，就会造成后续安全问题。

\textbf{零知识证明}（Zero-Knowledge Proof, ZKP）的目标是让证明者能够向验证者证明“自己确实知道某个秘密”，同时又不泄露这个秘密本身\cite{GMR1989}。对于本实验而言，Alice希望向Bob证明自己知道方程
\[
x^3+x+5=35
\]
的一个解，但又不希望把这个解直接告诉Bob（我知道答案但就是不告诉你）。题目中给出的已知正确解是 $x=3$，因此实验任务可以概括为：\textbf{在不公开 $x$ 的前提下完成证明生成与证明验证}。

\subsection{零知识证明的基本性质}
一个合格的零知识证明系统通常需要满足以下三个性质：

\begin{enumerate}
    \item \textbf{完备性（Completeness）}：若证明者确实知道秘密，并且诚实执行协议，则验证者应当接受证明。
    \item \textbf{可靠性（Soundness）}：若证明者并不知道秘密，则不应轻易伪造出可通过验证的证明。
    \item \textbf{零知识性（Zero-Knowledge）}：验证者在验证成功后，只知道“证明者掌握该秘密”，而不能从证明过程中恢复出秘密本身。
\end{enumerate}

因此，本实验的目标不是简单计算出 $x=3$，而是设计一个协议，使得Bob能够确认Alice知道解，但又不能直接得到解。

\subsection{从承诺开始理解证明原理}
在零知识证明中，\textbf{承诺（Commitment）} 是非常核心的思想。可以将其理解为：证明者先给出一个与秘密相关、但又不会直接泄露秘密的中间值，从而把自己当前的“说法”先固定下来。后续验证者再基于这个承诺提出挑战，证明者据此给出响应。这样就形成了零知识证明中经典的“承诺-挑战-响应”结构。

本实验采用\textbf{Schnorr型知识证明协议}\cite{Schnorr1991}来模拟这一过程。设公开参数为素数 $p$、子群阶 $q$ 和生成元 $g$，令Alice掌握的秘密 witness 为 $x$，对应的公开值为：
\[
pk=g^x \bmod p.
\]
其中 $pk$ 是公开陈述，$x$ 是秘密 witness。Bob可以知道 $pk$，但不能从中直接恢复出 $x$。

在交互式Schnorr协议中，证明过程可分为以下四步：

\begin{enumerate}
    \item \textbf{承诺阶段}：Alice随机选择一次性随机数 $r\in \mathbb{Z}_q$，计算
    \[
    R=g^r \bmod p.
    \]
    这里的 $R$ 就是承诺值。由于 $r$ 是随机的，因此 $R$ 不会直接泄露秘密 $x$。

    \item \textbf{挑战阶段}：Bob向Alice发送一个挑战值 $c$，要求Alice针对该挑战证明自己确实知道秘密。

    \item \textbf{响应阶段}：Alice根据随机数 $r$ 和秘密 $x$ 计算
    \[
    z=r+c\cdot x \bmod q.
    \]
    并把响应值 $z$ 发送给Bob。

    \item \textbf{验证阶段}：Bob检查如下等式是否成立：
    \[
    g^z \stackrel{?}{=} R\cdot pk^c \pmod p.
    \]
\end{enumerate}

上述等式成立的原因是：
\[
g^z=g^{r+cx}=g^r\cdot g^{cx}=R\cdot (g^x)^c=R\cdot pk^c \pmod p.
\]
因此，只要Alice确实知道秘密 $x$，她就能够针对挑战构造出满足等式的响应；而不知道 $x$ 的人则难以伪造通过验证的证明。

\subsection{非交互化处理}
标准Schnorr协议是交互式协议，需要Alice和Bob来回通信。为了便于程序实现，本实验采用\textbf{Fiat-Shamir变换}\cite{FiatShamir1986}，将原本由Bob给出的挑战改为通过哈希函数自动生成：
\[
c = H(out, pk, R)\bmod q.
\]
这样一来，Alice可以在本地一次性完成承诺、挑战和响应的计算，最终输出证明
\[
(out, pk, R, z).
\]
Bob拿到证明后，只需重新计算挑战值，再检查验证等式是否成立即可。这就把原本的交互式证明改造成了非交互证明。

\subsection{进阶实验原理：基于libsnark的通用零知识证明}
如果希望把本实验进一步升级成更严格的通用零知识证明，而不仅仅是教学化的 Schnorr 演示，那么可以使用 libsnark 库来完成基于 R1CS 的 zkSNARK 证明。此时证明的对象不再只是“我知道一个秘密数”，而是“我知道一个私密输入 $x$，它满足公开方程 $x^3+x+5=35$”。

\subsubsection{约束系统的构造}
为了让 libsnark 处理该方程，可以先引入两个中间变量：
\[
t_1=x\cdot x,\qquad t_2=t_1\cdot x.
\]
于是原方程可被拆成以下三条约束：
\[
\begin{cases}
x \cdot x = t_1 \\
t_1 \cdot x = t_2 \\
(t_2 + x + 5)\cdot 1 = out
\end{cases}
\]
其中公开输入为 $out=35$，私密 witness 为：
\[
x=3,\qquad t_1=9,\qquad t_2=27.
\]

\subsubsection{R1CS与zkSNARK流程}
libsnark 接受的核心输入之一是 R1CS（Rank-1 Constraint System，秩一约束系统）。对本实验来说，上述三条关系都可以写成秩一约束，随后交由 zkSNARK 系统处理。整个流程通常包括：
\begin{enumerate}
    \item 构造约束系统；
    \item 生成证明密钥和验证密钥；
    \item 输入公开输入与私密 witness，生成证明；
    \item 验证者仅使用公开输入和验证密钥完成验证。
\end{enumerate}
与前文的方法相比，这种方式证明的已经不是“知道某个秘密”，而是“知道一个满足公开约束系统的私密输入”，因此更符合通用零知识证明的标准形式\cite{Libsnark}。

% \subsection{实验中的模拟方法}
% 严格来说，若要直接对方程
% \[
% x^3+x+5=35
% \]
% 本身构造通用零知识证明，通常需要将该算术关系转化为约束系统，再交给更完整的zkSNARK框架处理。考虑到本实验规模较小，且目标是理解零知识证明的核心流程，因此本文采用教学化模拟方法：

% \begin{enumerate}
%     \item 先对公共方程进行约束检查，确认其整数解为 $x=3$。
%     \item 将该解视为Alice掌握的秘密 witness。
%     \item 再利用该 witness 构造 Schnorr/Fiat-Shamir 型非交互证明，模拟“知道解但不直接公开解”的过程。
% \end{enumerate}

% 这种实现方式虽然不是通用zkSNARK系统，但能够完整展示零知识证明中最核心的协议结构，即承诺、挑战、响应与验证。




\section{实验过程}
\subsection{实验环境}

本实验分为基础实现与进阶实现两个部分，分别采用不同的工具与库。

在基础实现中，主要使用 Python 标准库完成零知识证明流程的模拟，所使用主要库及其作用如下：
\begin{itemize}
    \item \texttt{hashlib}：用于实现 Fiat-Shamir 变换，通过哈希函数生成挑战值；
    \item \texttt{secrets}：用于生成密码学安全随机数，以模拟承诺阶段中的随机性；
\end{itemize}

在进阶实现中，采用 \texttt{libsnark} 库完成 zkSNARK 的完整流程。该库能够直接处理 R1CS（Rank-1 Constraint System）约束系统，并提供从密钥生成、证明生成到证明验证的一体化接口，极大简化了零知识证明系统的实现过程。

本实验中涉及的核心接口如表 \ref{tab:libsnark_api} 所示。

{
\renewcommand{\arraystretch}{1.3} % 行距拉伸

\begin{table}[h]
\centering
\small
\caption{libsnark 核心接口说明}
\label{tab:libsnark_api}

\rowcolors{2}{gray!10}{white} % 从第二行开始斑马纹

\begin{tabular}{>{\raggedright\arraybackslash}p{5.2cm} >{\raggedright\arraybackslash}p{8.5cm}}
\toprule
\rowcolor{gray!25}
\textbf{接口名称} & \textbf{功能说明} \\
\midrule

\texttt{protoboard<FieldT>} & 用于构建并维护整个约束系统，包含变量定义及约束关系。 \\

\texttt{pb\_variable<FieldT>} & 用于定义电路中的变量，包括公开输入（public input）和私密见证（witness）。 \\

\texttt{r1cs\_constraint<FieldT>} & 用于描述单条 R1CS 约束，即形如 $a \cdot b = c$ 的秩一约束关系。 \\

\texttt{r1cs\_gg\_ppzksnark\_generator} & 根据约束系统生成证明密钥（proving key）和验证密钥（verification key）。 \\

\texttt{r1cs\_gg\_ppzksnark\_prover} & 在给定 witness 的情况下，生成对应的零知识证明（proof）。 \\

\texttt{r1cs\_gg\_ppzksnark\_verifier\_strong\_IC} & 使用验证密钥与公开输入，对证明进行验证，判断其有效性。 \\

\bottomrule
\end{tabular}
\end{table}

}

\subsection{实验参数}
为了便于代码演示，本实验选用一个较小但足以说明过程的安全素数群：
\[
p=1019,\quad q=509,\quad g=3.
\]
其中 $q \mid (p-1)$，$g$ 是该子群的生成元。由于本实验只用于模拟而非真实部署，为了方便本地计算，因此没有采用大整数群参数。


\subsection{实验模拟步骤}
\subsubsection{步骤一：验证公共方程并确定 witness}
先对方程
\[
x^3+x+5=35
\]
进行枚举检查，确认在实验采用的整数范围内它只有一个整数解：
\[
x=3.
\]
这一步的作用是明确实验中的公共语句与秘密 witness。其中，方程和 $out=35$ 是公开信息，而 $x=3$ 只由Alice掌握。

\subsubsection{步骤二：生成公开陈述}
利用 witness 生成公开陈述：
\[
pk=g^x \bmod p = 3^3 \bmod 1019 = 27.
\]
在后续证明中，Alice并不直接发送 $x$，而是围绕公开值 $pk$ 构造零知识证明。

\subsubsection{步骤三：模拟承诺、挑战与响应}
程序随机生成一次性随机数 $r$，先计算承诺值 $R$，再通过哈希函数生成挑战值 $c$，最后计算响应值 $z$，最终得到证明对象：
\[
(out, pk, R, z).
\]
由于挑战由哈希函数自动生成，因此这里实现的是非交互版本，Alice可以一次性生成完整证明。

\subsubsection{步骤四：验证证明}
验证者拿到证明后，重新计算挑战值，并检查
\[
g^z = R\cdot pk^c \pmod p
\]
是否成立。若成立，则说明证明与Alice掌握的秘密一致；若不成立，则拒绝该证明。

\subsection{核心代码说明}
程序的核心由四个部分组成：

\begin{enumerate}
    \item \texttt{equation(x)}：计算 $x^3+x+5$。
    \item \texttt{find\_integer\_solutions(out)}：枚举公共方程的整数解。
    \item \texttt{generate\_proof(out, witness)}：生成Schnorr/Fiat-Shamir证明。
    \item \texttt{verify\_proof(proof)}：完成证明校验。
\end{enumerate}

生成证明和证明校验的关键实现代码如下。

\begin{lstlisting}[language=Python,caption={零知识证明核心实现},label={lst:zkp}]
import hashlib
import secrets
from dataclasses import dataclass

P = 1019
Q = 509
G = 3

def generate_proof(out: int, witness: int) -> Proof:
    nonce = secrets.randbelow(Q - 1) + 1
    public_key = pow(G, witness % Q, P)
    commitment = pow(G, nonce, P)
    challenge = hash_to_scalar(out, public_key, commitment)
    response = (nonce + challenge * (witness % Q)) % Q
    return Proof(out, public_key, commitment, response)

def verify_proof(proof: Proof) -> bool:
    challenge = hash_to_scalar(proof.out, proof.public_key, proof.commitment)
    left = pow(G, proof.response, P)
    right = (proof.commitment * pow(proof.public_key, challenge, P)) % P
    return left == right
\end{lstlisting}

\newpage

\subsection{程序运行结果}
执行程序输出如下。由于随机数每次都会变化，因此承诺值与响应值会有所不同，但验证结论保持不变。

\begin{lstlisting}[caption={程序运行输出}]
=== 公共语句 ===
方程: x^3 + x + 5 = 35
公开参数: p=1019, q=509, g=3

=== witness 检查 ===
整数范围[-100, 100]内的解: [3]
Alice 持有的 witness x = 3
代入验证: 3^3 + 3 + 5 = 35

=== 证明生成 ===
{
  "out": 35,
  "public_key": 27,
  "commitment": 836,
  "response": 342
}

=== 证明验证 ===
验证结果: True

=== 篡改测试 ===
篡改后的验证结果: False
\end{lstlisting}

从结果可以看出：
\begin{enumerate}
    \item 正常生成的证明能够通过验证。
    \item 只要对证明中的响应值做微小篡改，验证就会失败。
    \item 这说明程序已经实现了完整的“证明生成-证明验证”闭环。
\end{enumerate}

\subsection{进阶实验：libsnark实现模拟过程}
\subsubsection{步骤一：构造R1CS约束系统}
在进阶实现中，不再直接围绕 Schnorr 承诺构造证明，而是先把方程
\[
x^3+x+5=35
\]
拆成约束系统。程序中通过 \texttt{protoboard} 定义变量 \texttt{out}、\texttt{x}、\texttt{x2}、\texttt{x3}，再依次加入三条约束：
\[
x\cdot x=x2,\qquad x2\cdot x=x3,\qquad (x3+x+5)\cdot 1=out.
\]

\subsubsection{步骤二：设置公开输入与私密 witness}
这里取公开输入 \texttt{out=35}，私密 witness 为：
\[
x=3,\qquad x2=9,\qquad x3=27.
\]
如果这些值满足全部约束，则说明 witness 与公开方程是一致的。

\subsubsection{步骤三：生成证明密钥、验证密钥与证明}
在 libsnark 中，先调用
\texttt{r1cs\_gg\_ppzksnark\_generator}
根据约束系统生成密钥对，再调用
\texttt{r1cs\_gg\_ppzksnark\_prover}
输入公开输入和私密 witness 生成证明。

\subsubsection{步骤四：验证证明}
验证阶段调用
\texttt{r1cs\_gg\_ppzksnark\_verifier\_strong\_IC}，
只输入公开输入和证明进行校验。若验证成功，则说明证明者确实知道一个满足约束系统的私密输入，而无需暴露该输入本身。

\subsubsection{进阶代码说明}
进阶实验核心流程如下：

\begin{lstlisting}[language=C++,caption={libsnark进阶实现核心流程}]
typedef alt_bn128_pp ppT;
typedef Fr<ppT> FieldT;

ppT::init_public_params();

protoboard<FieldT> pb;
pb_variable<FieldT> out, x, x2, x3;

out.allocate(pb, "out");
x.allocate(pb, "x");
x2.allocate(pb, "x2");
x3.allocate(pb, "x3");
pb.set_input_sizes(1);

pb.add_r1cs_constraint(r1cs_constraint<FieldT>(x, x, x2));
pb.add_r1cs_constraint(r1cs_constraint<FieldT>(x2, x, x3));
pb.add_r1cs_constraint(
    r1cs_constraint<FieldT>(x3 + x + FieldT("5"), FieldT::one(), out)
);

pb.val(out) = FieldT("35");
pb.val(x) = FieldT("3");
pb.val(x2) = FieldT("9");
pb.val(x3) = FieldT("27");

auto cs = pb.get_constraint_system();
auto keypair = r1cs_gg_ppzksnark_generator<ppT>(cs);
auto proof = r1cs_gg_ppzksnark_prover<ppT>(
    keypair.pk, pb.primary_input(), pb.auxiliary_input()
);
bool ok = r1cs_gg_ppzksnark_verifier_strong_IC<ppT>(
    keypair.vk, pb.primary_input(), proof
);
\end{lstlisting}

\section{零知识证明-可靠性探究}
实验做完了，我们来探究一下这个0知识证明为啥可靠呢，从上课几周的经验来看，现在的加密可靠性都依赖于巨大计算量：从协议执行角度看，Schnorr/Fiat-Shamir证明本身的计算量并不高。证明生成阶段主要涉及两次模幂运算：一次计算承诺值 $R=g^r \bmod p$，一次计算公开值 $pk=g^x \bmod p$；验证阶段同样主要涉及两次模幂运算：一次计算 $g^z \bmod p$，一次计算 $pk^c \bmod p$，再配合一次模乘即可完成校验。因此，对诚实参与方而言，该协议的运行开销大致是常数次模幂运算，复杂度可以近似看作 $O(\log q)$ 次乘法级别。

从攻击者角度看，真正决定可靠性的，是恢复秘密 $x$ 或伪造响应所需要付出的计算代价。在本实验参数下，$q=509$，秘密指数的取值空间最多只有 509 种。也就是说，如果攻击者直接暴力枚举 $x=0,1,\dots,508$，逐个检查
\[
g^x \stackrel{?}{=} pk \pmod p
\]
是否成立，最多只需要 509 次尝试就可以恢复秘密。对于现代计算机而言，这个代价几乎可以忽略，因此这里的小素数参数只能用于教学演示，而不能提供真实安全性。

若把参数提升到实际密码系统常见的大素数规模，例如令子群阶 $q$ 约为 $2^{256}$，那么暴力枚举秘密的复杂度将提升到
\[
2^{256}\approx 1.16\times 10^{77}
\]
次尝试。这一数量级已经远远超过现实中任何计算设备能够承受的范围。即使假设一台设备每秒可以完成 $10^{12}$ 次尝试，穷举全部空间所需时间仍然约为
\[
\frac{2^{256}}{10^{12}} \approx 1.16\times 10^{65}\ \text{秒},
\]
折合约为
\[
3.7\times 10^{57}\ \text{年},
\]
这在实际上是不可完成的。

因此，本方法的可靠性可以这样理解：协议本身对诚实双方来说只需要少量模幂运算，开销较低；而对攻击者而言，若参数足够大，恢复秘密或伪造证明的代价会随着子群阶 $q$ 的增大呈指数级上升。也正因为如此，实验中选取的小素数适合展示协议流程，而真正部署时必须使用大素数参数，才能使该方法具备实际安全性。

\section{实验总结 \& 心得体会}
本次实验围绕“证明自己知道方程解”这一目标，完成了一个简化的零知识证明原型。通过实验可以看出，零知识证明并不只是一个抽象概念，而是可以被拆解为明确的协议步骤：先做承诺，再生成挑战，随后给出响应，最后由验证方检查等式。

从实现角度看，本实验采用了Schnorr协议结合Fiat-Shamir变换的非交互形式。这样做的好处是：代码量小、逻辑清晰、能够完整覆盖证明生成和证明验证两个关键环节。虽然它还不是一个通用的zkSNARK实现，但已经足以帮助理解零知识证明的基本思想与工程结构。

但是存在一个巨大槽点，Bob只能知道你手里有一个答案，怎么确实是真答案？，所以我还尝试了进阶实验：把本实验扩展成“对任意多项式约束做真正的零知识证明”。尝试过程中接触了新的库、新的流程，收获满满。


% ======================== 参考文献模块（调用bib文件） ========================
\newpage % 参考文献单独分页
\bibliographystyle{plain} % 参考文献样式（plain/unsrt/gb7714-2015等）
\renewcommand{\refname}{\Large\textbf{参考文献}} % 标题居中加粗
\bibliography{refs} % 调用refs.bib文件（无需写后缀.bib）
% ======================== 参考文献模块结束 ========================
\end{document}
